\begin{abstract}
We propose a constructive foundation for set theory in which the primitive
notion is not membership but \emph{discrimination}: the act of
interrogating a collection via two independent binary predicates
($\tau$ and $\sigma$), each valued in $\GF(2)$. Each predicate
independently partitions the collection into
\emph{comprehended} and \emph{uncomprehended}; the four-cell
membership profile is the product of these two independent
partitions.
A growing discrimination structure $\kappa$ records all discriminations
performed. Sets are equivalence classes under $\kappa$, membership is a
ternary relation (element, collection, discriminator), and the Boolean
law of excluded middle is not assumed globally but
constructed locally on dimensions where the algebra and
regulatory axioms jointly support it. We show that extensionality, a restricted
comprehension principle, and a form of foundation emerge as consequences
of the algebra together with two regulatory axioms
(profile linearity (Definition \ref{def:profile-linear}) and the operationalist principle),
rather than as independent axioms. In particular, Russell's predicate
is excluded by the profile linearity axiom: the Russell predicate
produces an affine $\tau$-profile, which is inadmissible for a
linear discriminator. The restriction is relocated from set
formation (as in ZFC) to discriminator admissibility.
\end{abstract}

